\documentclass[a4paper]{article}

\usepackage{graphicx}
\usepackage{amsmath,amssymb}
\usepackage{minted}
\usepackage{multirow}
\usepackage{float}
\usepackage{todonotes}
\usepackage{forest}
\usepackage{xstring}
\usepackage{xcolor}
\usepackage[most]{tcolorbox}
\usepackage{xparse}
\usepackage{natbib}

\usetikzlibrary{graphs,graphdrawing}
\usegdlibrary{layered}

% for external graphviz
\input{/home/donkarlo/Dropbox/repo/nd_language_ai_project/src/nd_language_ai/formal/computer/programming/tex/latex/code_clips/external_graphviz.tex}

% for tags
\input{/home/donkarlo/Dropbox/repo/nd_language_ai_project/src/nd_language_ai/formal/computer/programming/tex/latex/parts/control_sequence/kind/semtag/semtag.tex}

% for links
\input{/home/donkarlo/Dropbox/repo/nd_language_ai_project/src/nd_language_ai/formal/computer/programming/tex/latex/code_clips/seehere.tex}

\title{}
\date{}

\begin{document}
    \maketitle
    
    Man benutzt es mit einem anderen Verb.
    Es bleibt ein Verb.
    
    Beispiel:
    Ich versuche, Deutsch zu lernen.
    I try to learn German.
    
    \bigskip
    
    \textbf{das + Verb:}
    
    Man macht aus dem Verb ein Nomen.
    Es ist ein Ding oder eine Aktivität.
    
    Beispiel:
    Das Lernen ist wichtig.
    Learning is important.
    
    \bigskip
    
    \textbf{Kurz:}
    
    zu lernen = Verb
    
    das Lernen = Nomen
    
    
    \section{das + Verb (Infinitiv, großgeschrieben)}
        Das Rauchen ist verboten: Smoking is forbidden.
        
        Ich liebe das Reisen: I love traveling.
    
    
    \section{Infinitiv mit „zu“}
        Ich versuche, Deutsch zu lernen: I try to learn German.
        
        Es ist wichtig, genug zu schlafen: It is important to sleep enough.
    %\bibliographystyle{plainnat}
    %\bibliography{/home/donkarlo/Dropbox/repo/nd_shared_research/bibliography/bibliography.bib}
\end{document}